\documentclass[12pt,a4paper]{article}
\usepackage[utf8]{inputenc}
\usepackage[T1]{fontenc}
\usepackage[spanish]{babel}
\usepackage{geometry}
\usepackage{xcolor}
\usepackage{listings}
\usepackage{booktabs}
\usepackage{longtable}
\usepackage{hyperref}
\usepackage{fancyhdr}
\usepackage{titlesec}
\usepackage{amsmath}
\usepackage{enumitem}
\usepackage{float}
\usepackage{array}
\usepackage{tcolorbox}

\geometry{margin=2.5cm}

\definecolor{primary}{RGB}{31,73,125}
\definecolor{secondary}{RGB}{68,114,196}
\definecolor{codebg}{RGB}{245,245,245}
\definecolor{codeframe}{RGB}{200,200,200}
\definecolor{warningbg}{RGB}{255,243,205}
\definecolor{successbg}{RGB}{212,237,218}
\definecolor{errorbg}{RGB}{248,215,218}

\lstset{
    backgroundcolor=\color{codebg},
    frame=single,
    rulecolor=\color{codeframe},
    basicstyle=\ttfamily\small,
    keywordstyle=\color{primary}\bfseries,
    commentstyle=\color{gray}\itshape,
    stringstyle=\color{secondary},
    breaklines=true,
    showstringspaces=false,
    numbers=left,
    numberstyle=\tiny\color{gray},
    language=Python,
}

\pagestyle{fancy}
\fancyhf{}
\fancyhead[L]{\textcolor{primary}{\textbf{Manual Técnico}}}
\fancyhead[R]{\textcolor{primary}{GeoMIP — k-Particiones}}
\fancyfoot[C]{\thepage}
\fancyfoot[L]{\small Universidad de Caldas}
\fancyfoot[R]{\small 2026}

\titleformat{\section}{\Large\bfseries\color{primary}}{}{0em}{\thesection.\quad}
\titleformat{\subsection}{\large\bfseries\color{secondary}}{}{0em}{\thesubsection.\quad}
\titleformat{\subsubsection}{\normalsize\bfseries}{}{0em}{\thesubsubsection.\quad}

\hypersetup{
    colorlinks=true,
    linkcolor=primary,
    urlcolor=secondary,
}

\tcbuselibrary{skins}
\newtcolorbox{notebox}[1][]{
    colback=warningbg, colframe=orange!70!black,
    title={\textbf{Nota}}, fonttitle=\small, #1
}
\newtcolorbox{bugbox}[1][]{
    colback=errorbg, colframe=red!60!black,
    title={\textbf{Bug corregido}}, fonttitle=\small, #1
}
\newtcolorbox{optbox}[1][]{
    colback=successbg, colframe=green!50!black,
    title={\textbf{Optimización}}, fonttitle=\small, #1
}

\begin{document}

% ─────────────────────────── PORTADA ───────────────────────────
\begin{titlepage}
    \centering
    \vspace*{1cm}
    {\Huge\bfseries\color{primary} Manual Técnico\\[0.5em]
    \Large Sistema GeoMIP\\[0.3em]
    \large Algoritmo k-MIP Geométrico}
    \vspace{1.5cm}

    \rule{\linewidth}{2pt}
    \vspace{0.5cm}

    {\large\textbf{Proyecto de Análisis y Diseño de Algoritmos}}\\[0.3em]
    {\large Ingeniería en Sistemas y Computación}\\[0.3em]
    {\large Facultad de Ingeniería — Universidad de Caldas}
    \vspace{1cm}

    \rule{\linewidth}{0.5pt}
    \vspace{0.8cm}

    {\large\textbf{Autores:}}\\[0.5em]
    {\large Jorge Ivan Garcia Torres}\\[0.3em]
    {\large Camilo Botero Merchor}
    \vspace{1cm}

    \rule{\linewidth}{0.5pt}
    \vspace{0.5cm}
    {\large Semestre 2026-1}

    \vfill
    {\small\color{gray} proyecto-analisis-kpartitions / GeoMIP}
\end{titlepage}

\tableofcontents
\newpage

% ─────────────────────────── 1. INSTALACIÓN ───────────────────────────
\section{Instalación y Configuración del Entorno}

\subsection{Requisitos del sistema}
\begin{longtable}{ll}
    \toprule
    Componente & Requisito \\
    \midrule
    \endhead
    Sistema operativo & Windows 10/11, Linux o macOS \\
    Python & 3.11 o superior \\
    RAM mínima & 4 GB (8 GB recomendado para N22A, N25A) \\
    Disco & 2 GB libres (N25A.csv $\approx$ 800 MB) \\
    \bottomrule
\end{longtable}

\subsection{Creación del entorno virtual}
\begin{lstlisting}[language=bash]
# Desde la raiz del repositorio
python -m venv .venv

# Windows (Git Bash)
source .venv/Scripts/activate

# Linux / macOS
source .venv/bin/activate
\end{lstlisting}

\subsection{Instalación de dependencias}
\begin{lstlisting}[language=bash]
pip install numpy openpyxl pandas pyinstrument
\end{lstlisting}

\subsection{Verificación}
\begin{lstlisting}[language=bash]
python -c "import numpy, openpyxl, pandas; print('OK')"
\end{lstlisting}

% ─────────────────────────── 2. MÓDULOS ───────────────────────────
\section{Referencia de Módulos}

\subsection{generar\_tpms.py}

\textbf{Propósito:} Genera los archivos CSV de matrices de transición de probabilidad (TPM).

\textbf{Función principal:}
\begin{lstlisting}
def generar_tpm(n_nodos: int, sufijo: str, carpeta: Path,
                semilla: int = 0, forzar: bool = False) -> Path
\end{lstlisting}

\begin{longtable}{lll}
    \toprule
    Parámetro & Tipo & Descripción \\
    \midrule
    \endhead
    \texttt{n\_nodos} & int & Número de nodos del sistema (10, 15, 20, 22, 25) \\
    \texttt{sufijo}  & str & Letra identificadora del dataset ('A', 'B', ...) \\
    \texttt{carpeta} & Path & Ruta destino del archivo CSV \\
    \texttt{semilla} & int & Semilla numpy (default: 0) \\
    \texttt{forzar}  & bool & Si True, regenera aunque ya exista \\
    \bottomrule
\end{longtable}

\textbf{Estructura de la TPM generada:}
\begin{lstlisting}
# Matriz 2^N x N de enteros 0/1
np.random.seed(semilla)
matriz = np.random.randint(2, size=(2**n_nodos, n_nodos), dtype=np.uint8)
np.savetxt(ruta, matriz, delimiter=",", fmt="%d")
\end{lstlisting}

\begin{notebox}
La semilla numpy garantiza reproducibilidad. Con \texttt{semilla=0} los archivos son idénticos en cualquier máquina. Cambiar la semilla genera datasets distintos.
\end{notebox}

\subsection{excel\_loader.py}

\textbf{Función principal:}
\begin{lstlisting}
def cargar_datos_excel(ruta: str) -> list[dict]
\end{lstlisting}

\textbf{Estructura del dict devuelto por caso:}
\begin{lstlisting}
{
    "hoja":           str,   # ej. "10A-Elementos"
    "fila_excel":     int,   # numero de fila en Excel (desde 6)
    "prueba":         int,   # numero de prueba
    "estado_inicial": str,   # binario, ej. "1000000000"
    "condicion":      str,   # letras, ej. "ABCDEFGHIJ"
    "subsistema":     str,   # letras (alcance), ej. "ABCDEFGHI"
    "mecanismo":      str,   # letras, ej. "ACEGI"
}
\end{lstlisting}

\begin{notebox}
El estado inicial se sanitiza automáticamente: ``1,000,000,000'' → ``1000000000''. El corte anticipado tras 15 filas vacías consecutivas evita iterar las 1003 filas vacías que tiene cada hoja.
\end{notebox}

\subsection{excel\_writer.py}

\textbf{Mapeo de columnas (COLUMNAS\_RESULTADOS):}
\begin{lstlisting}
COLUMNAS_RESULTADOS = {
    (2, "QNodes"):    ("D", "E", "F"),
    (2, "Geometric"): ("G", "H", "I"),
    (3, "QNodes"):    ("J", "K", "L"),
    (3, "Geometric"): ("M", "N", "O"),
    (4, "QNodes"):    ("P", "Q", "R"),
    (4, "Geometric"): ("S", "T", "U"),
    (5, "QNodes"):    ("V", "W", "X"),
    (5, "Geometric"): ("Y", "Z", "AA"),
}
\end{lstlisting}

\textbf{Función de escritura:}
\begin{lstlisting}
def escribir_resultado_ws(ws, fila: int, k: int, estrategia: str,
                          particion: str, perdida: float,
                          tiempo: float) -> None
\end{lstlisting}

\subsection{KPartitionGeometricSIA (geometric\_k.py)}

\subsubsection{Constructor}
\begin{lstlisting}
def __init__(self, gestor: Manager, k: int = 2)
\end{lstlisting}

\begin{longtable}{lll}
    \toprule
    Atributo & Tipo & Descripción \\
    \midrule
    \endhead
    \texttt{k} & int & Número de particiones (2--5) \\
    \texttt{tabla\_transiciones} & dict & Tabla T: (estado\_i, estado\_j) → [costos] \\
    \texttt{\_flat\_data} & list[ndarray] & TPM aplanada por nodo para acceso O(1) \\
    \texttt{\_caminos} & dict & BFS: nivel → lista de estados \\
    \texttt{\_estado\_int\_cache} & dict & Cache de conversión tuple → int \\
    \texttt{memoria\_particiones} & dict & key → (emd, distribución) \\
    \bottomrule
\end{longtable}

\subsubsection{Método principal}
\begin{lstlisting}
@profile(context={TYPE_TAG: GEOMETRIC_ANALYSIS_TAG})
def aplicar_estrategia(self, condicion: str, alcance: str,
                       mecanismo: str, tpm: np.ndarray) -> Solution
\end{lstlisting}

\begin{notebox}
\textbf{Orden de argumentos:} el runner prueba dos órdenes automáticamente: \\
\texttt{(condicion, mecanismo, alcance, tpm)} y \texttt{(condicion, alcance, mecanismo, tpm)}.
\end{notebox}

\subsubsection{\_construir\_tabla}
BFS desde estado inicial hasta estado final (Hamming creciente). Complejidad: $O(N \cdot 2^N)$ operaciones de costo.

\begin{lstlisting}
def _construir_tabla(self) -> None:
    # nivel 0: [estado_inicial]
    # nivel 1: vecinos a distancia Hamming 1
    # ...
    # nivel N: [estado_final]
\end{lstlisting}

\subsubsection{\_candidatos\_k2 (OPTIMIZADO)}

\begin{optbox}
Para k=2, en vez de generar $S(n,2) = 2^{n-1}-1$ candidatos, se generan $\sim 1.5 \cdot n$:
\begin{itemize}
    \item \textbf{Tipo A:} separar cada variable futura individualmente ($n$ candidatos).
    \item \textbf{Tipo B:} selección BFS por nivel hasta la mitad ($\sim n/2$ candidatos).
\end{itemize}
Para $n=10$: 14 candidatos vs 511 del original ($\sim 34\times$ más rápido).
\end{optbox}

\begin{bugbox}
\textbf{Índices desplazados:} los futuros usan índices $[n\_pres, \ldots, n\_pres + n\_vars - 1]$ y los presentes $[0, \ldots, n\_pres-1]$. Sin este desplazamiento, \texttt{\_evaluar\_candidato} confundía futuros con presentes cuando $n\_vars \leq n\_pres$, produciendo el error \texttt{shapes (225,) (15,) could not be broadcast}.
\end{bugbox}

\subsubsection{\_evaluar\_candidato}
\begin{lstlisting}
def _evaluar_candidato(self, grupos_vars: list[list[int]]
                       ) -> tuple[float, np.ndarray] | None:
    # Para cada grupo:
    #   - v < n_pres  => presente: dims_ncubos[v]
    #   - v >= n_pres => futuro:   indices_ncubos[v - n_pres]
    # Calcula biparticion y distribucion marginal
    # Para k>2: producto tensorial (np.kron)
    # Verifica compatibilidad de shapes antes de EMD
    # Retorna (emd, dist) o None si hay error
\end{lstlisting}

\subsubsection{\_find\_k\_mip}
\begin{lstlisting}
def _find_k_mip(self) -> tuple:
    self._construir_tabla()
    candidatos = self._candidatos_geometricos(self.k)
    for grupos in candidatos:
        resultado = self._evaluar_candidato(grupos)
        if resultado is not None:
            self.memoria_particiones[key] = resultado
            if resultado[0] == 0.0:   # Early-exit
                return key
    return min(self.memoria_particiones,
               key=lambda k: self.memoria_particiones[k][0])
\end{lstlisting}

\subsection{fill\_excel\_runner.py}

\subsubsection{\_convertir\_a\_binario}
\begin{lstlisting}
def _convertir_a_binario(texto: str, n_bits: int) -> str:
    # Si ya es binario (solo 0/1), devuelve tal cual ajustado a n_bits
    # Si son letras: A=bit0, B=bit1, ..., T=bit19
    # Ej: "ACEGI", n=10 -> "1010101010"
    #     "BCDEFGHIJ", n=10 -> "0111111111"
\end{lstlisting}

\subsubsection{Timeout por proceso}
\begin{lstlisting}
def _ejecutar_con_timeout(target, args: tuple, timeout: int) -> dict:
    q = multiprocessing.Queue()
    p = multiprocessing.Process(target=target, args=(q,) + args)
    p.start()
    p.join(timeout=timeout)
    if p.is_alive():
        p.terminate()
        p.join()
        return {"ok": False, "error": f"Timeout ({timeout}s)"}
    return q.get() if not q.empty() else {"ok": False, "error": "Sin resultado"}
\end{lstlisting}

% ─────────────────────────── 3. ERRORES CONOCIDOS ───────────────────────────
\section{Errores Conocidos y Soluciones}

\subsection{SyntaxError: unmatched ')'}
\textbf{Causa:} Corrupción del archivo en línea 4 de \texttt{fill\_excel\_runner.py}. La línea \texttt{fromfile\_\_).resolve()from pathlib import Path} es una fusión de dos líneas.

\textbf{Solución:}
\begin{lstlisting}
# Reemplazar la linea corrupta por:
from pathlib import Path
THIS_FILE = Path(__file__).resolve()
\end{lstlisting}

\subsection{SyntaxError en excel\_writer.py línea 6}
\textbf{Causa:} Corrupción similar. La línea \texttt{\}' en el archivo Excel.")    (2, "QNodes"): ("D", "E", "F"),} fusiona el cierre del raise con la primera entrada del diccionario.

\textbf{Solución:} Separar en las dos líneas correspondientes y restaurar el orden lógico del archivo.

\subsection{operands could not be broadcast (225,) vs (15,)}
\textbf{Causa:} Índices de futuros sin desplazamiento por \texttt{n\_pres} en los candidatos.

\textbf{Solución:} En \texttt{\_candidatos\_k2} y \texttt{\_candidatos\_greedy\_kn}, todos los futuros deben ser \texttt{n\_pres + i} en vez de \texttt{i}.

\subsection{KeyError: 'Worksheet 25A-Elementos does not exist'}
\textbf{Causa:} La hoja tiene espacio al final (\texttt{'25A-Elementos '}) en el workbook.

\textbf{Solución:}
\begin{lstlisting}
hoja_real = next((s for s in wb.sheetnames if s.strip() == hoja), hoja)
ws = wb[hoja_real]
\end{lstlisting}

\subsection{FileNotFoundError: N20A.csv / N22A.csv / N25A.csv}
\textbf{Causa:} Los archivos TPM no fueron generados (no se incluyen en el repositorio por tamaño).

\textbf{Solución:}
\begin{lstlisting}[language=bash]
python generar_tpms.py --solo 20 22 25
\end{lstlisting}

\subsection{pyinstrument: session has 279756 samples}
\textbf{Causa:} El decorador \texttt{@profile} de pyinstrument acumula demasiadas muestras en ejecuciones largas (N=15+), causando overhead y mensajes de advertencia.

\textbf{Solución:} El nuevo runner ejecuta cada caso en un proceso hijo aislado (\texttt{multiprocessing}), evitando que el profiler del padre interfiera.

% ─────────────────────────── 4. MÉTRICAS ───────────────────────────
\section{Métricas del Benchmark}

\subsection{Definiciones (Tabla 5.1)}

\textbf{Tasa de acierto exacto:}
\[ \text{Tasa} = \frac{\text{casos con } J(\pi_{ref}, \pi_{geo}) < 10^{-6}}{N_{total}} \times 100\% \]

\textbf{Error relativo en } $\Phi$:
\[ E_{rel} = \frac{|\Phi_{ref} - \Phi_{geo}|}{\Phi_{ref}}, \quad E_{rel} = 0 \text{ si } \Phi_{ref} = 0 \]

\textbf{Distancia de Jaccard entre particiones:}
\[ J(\pi_a, \pi_b) = \min\left( \overline{J}_{dir}, \overline{J}_{cruz} \right) \]
donde $\overline{J}_{dir}$ y $\overline{J}_{cruz}$ son los promedios de las distancias de Jaccard en las dos posibles alineaciones de grupos.

\textbf{Speedup:}
\[ S_{rel} = \frac{T_{QNodes}}{T_{Geometric}} \]

\subsection{Clasificación de calidad}
\begin{longtable}{lrrr}
    \toprule
    Clasificación & Tasa de acierto & Error relativo & Jaccard \\
    \midrule
    \endhead
    Excelente    & $> 90\%$ & $< 1\%$  & $< 0.1$ \\
    Bueno        & $> 80\%$ & $< 5\%$  & $< 0.2$ \\
    Aceptable    & $> 70\%$ & $< 10\%$ & $< 0.3$ \\
    Insuficiente & $\leq 70\%$ & $\geq 10\%$ & $\geq 0.3$ \\
    \bottomrule
\end{longtable}

% ─────────────────────────── 5. GUÍA DE OPTIMIZACIÓN ───────────────────────────
\section{Guía de Optimización Futura}

\subsection{Optimizaciones implementadas}
\begin{enumerate}
    \item \textbf{Reducción de candidatos k=2:} $S(n,2) \to 1.5n$ ($\sim 34\times$ en n=10).
    \item \textbf{Vectorización numpy en cálculo de costos:} diferencias como operación de array.
    \item \textbf{Cache de enteros de estado:} \texttt{\_estado\_int\_cache} evita reconvertir tuplas.
    \item \textbf{Early-exit en $\Phi=0$:} retorna inmediatamente si la primera partición es perfecta.
    \item \textbf{Timeout por proceso:} evita bloqueos en casos pesados.
    \item \textbf{Corte anticipado en excel\_loader:} 15 vacíos consecutivos → stop (antes: 1003 filas).
\end{enumerate}

\subsection{Optimizaciones pendientes}
\begin{enumerate}
    \item \textbf{Paralelismo en evaluación de candidatos:} cada EMD es independiente.
\begin{lstlisting}
from concurrent.futures import ProcessPoolExecutor
with ProcessPoolExecutor() as ex:
    resultados = list(ex.map(_evaluar_candidato, candidatos))
\end{lstlisting}
    \item \textbf{Cache de distribuciones marginales:} si el mismo (alcance, mecanismo) se repite entre diferentes k, reusar la distribución calculada.
    \item \textbf{Lazy loading de TPM:} cargar solo las filas necesarias para el estado inicial (para N=25 el archivo es de 800 MB pero solo se acceden $\sim 2$ filas por estado).
\end{enumerate}

\end{document}